% Transfer Learning in RL
% Requires: \usepackage{booktabs}

\section{Shared Experimental Setup}

Task: \texttt{PointMaze\_UMaze-v3}, sparse reward, 300-step episodes,
\texttt{Discrete(9)} bang-bang actions. Two transfer directions are studied
under identical conditions:

\begin{itemize}
  \item \textbf{Study A} --- PPO pretrained on the corrupted environment, GFlowNet fine-tuned on the original.
  \item \textbf{Study B} --- GFlowNet pretrained on the corrupted environment, PPO fine-tuned on the original.
\end{itemize}

Each study runs three arms, $n=30$ seeds each, 300,000 environment steps per
phase:

\begin{itemize}
  \item \textbf{Pretrain} --- random init, on the \emph{corrupted} environment.
  \item \textbf{Fine-tuned} --- initialised by exact parameter copy of the pretrained policy, on the \emph{original} environment.
  \item \textbf{Scratch} (control) --- random init, on the \emph{original} environment.
\end{itemize}

Fine-tuned and scratch share algorithm, environment, budget and
hyperparameters, and differ only in initialisation; their contrast is the
primary outcome. Policies are evaluated frozen and deterministically on 100
fixed held-out instances. \texttt{hyperparameters.py} is byte-identical across
both studies (SHA-256 \texttt{d01ab117\dots}).

Three corruptions are used:

\begin{itemize}
  \item \texttt{large\_multigoal} --- a $12\times9$ maze with 7 goal cells replaces the $5\times5$ UMaze.
  \item \texttt{action\_gain\_jitter} --- actuator gain $g\sim\mathcal{U}(0.9,1.1)$, drawn per episode, unobservable. Since $g=1$ is interior to the support, this is domain randomisation, not a conflicting task.
  \item \texttt{negate\_both} --- $u\mapsto-u$; prior work, included as a reference and not re-run.
\end{itemize}

The transfer is an exact parameter copy, verified on all 30 seeds of every
condition (maximum absolute logit deviation $0.00$ over 512 probe observations;
greedy-action agreement $1.000$). One head cannot transfer in either direction
--- the GFlowNet's $\log Z$ head in Study A, PPO's critic in Study B --- and is
randomly initialised in both the fine-tuned and the scratch arm, so it is never
a difference between them. The scratch control is bit-identical across all
corruptions within a study (maximum discrepancy $0$), so corruptions are
compared against a common baseline.

\clearpage
\section{Study A: PPO to GFlowNet}

PPO pretrains on the corrupted environment and a GFlowNet is fine-tuned on the
original. The control is a GFlowNet trained from scratch, so the contrast holds
the fine-tuning algorithm fixed.

\begin{table}[htbp]
\centering
\setlength{\tabcolsep}{4pt}
\caption{Study A primary contrast: fine-tuned $-$ scratch on the original
environment. Both arms are GFlowNet. Positive favours fine-tuning, except
steps-to-goal where lower is better. Hedges' $g$ in parentheses;
${}^{*}p<0.05$, ${}^{**}p<0.01$, Holm-corrected over the four metrics.}
\small
\begin{tabular}{lcccc}
\toprule
Corruption & Success rate & Mean return & Steps to goal & Path efficiency \\
\midrule
\texttt{large\_multigoal} & $-0.007$ & $+11.1$ & $-11.5^{**}$ & $+0.021$ \\
 & \footnotesize$(-0.05)$ & \footnotesize$(0.57)$ & \footnotesize$(-0.92)$ & \footnotesize$(0.44)$ \\[3pt]
\texttt{action\_gain\_jitter} & $-0.021$ & $+16.8^{*}$ & $-7.4$ & $+0.013$ \\
 & \footnotesize$(-0.13)$ & \footnotesize$(0.73)$ & \footnotesize$(-0.59)$ & \footnotesize$(0.29)$ \\[3pt]
\texttt{negate\_both} (ref.) & $-0.025$ & $+4.0$ & $-7.2$ & $+0.004$ \\
 & \footnotesize$(-0.17)$ & \footnotesize$(0.22)$ & \footnotesize$(-0.55)$ & \footnotesize$(0.08)$ \\
\bottomrule
\end{tabular}
\end{table}

\begin{table}[htbp]
\centering
\setlength{\tabcolsep}{4pt}
\caption{Study A, final performance by arm, mean $\pm$ SD over $n=30$ seeds.
\emph{Native} is the pretrained policy on the corrupted environment it trained
on; \emph{zero-shot} is the same policy on the original environment without
adaptation. The GFlowNet scratch arm is shared across both corruptions.}
\small
\begin{tabular}{llcccc}
\toprule
Corruption & Arm & Success & Return & Steps & Path eff. \\
\midrule
\multicolumn{6}{l}{\texttt{large\_multigoal}} \\
 & Pretrain (native)    & $0.229 \pm 0.094$ & $16.8 \pm 8.5$ & $145.7 \pm 33.0$ & $0.851 \pm 0.092$ \\
 & Pretrain (zero-shot) & $0.274 \pm 0.122$ & $22.7 \pm 11.8$ & $92.4 \pm 24.6$ & $0.690 \pm 0.085$ \\
 & Fine-tuned           & $0.719 \pm 0.144$ & $62.3 \pm 20.5$ & $73.9 \pm 12.4$ & $0.701 \pm 0.050$ \\
\midrule
\multicolumn{6}{l}{\texttt{action\_gain\_jitter}} \\
 & Pretrain (native)    & $0.818 \pm 0.032$ & $183.5 \pm 12.1$ & $67.8 \pm 6.7$ & $0.898 \pm 0.024$ \\
 & Pretrain (zero-shot) & $0.821 \pm 0.030$ & $184.1 \pm 11.7$ & $68.2 \pm 6.0$ & $0.897 \pm 0.024$ \\
 & Fine-tuned           & $0.704 \pm 0.181$ & $68.0 \pm 26.7$ & $78.1 \pm 12.5$ & $0.693 \pm 0.044$ \\
\midrule
\multicolumn{2}{l}{GFlowNet scratch (control)} & $0.726 \pm 0.134$ & $51.2 \pm 18.1$ & $85.4 \pm 12.3$ & $0.679 \pm 0.046$ \\
\bottomrule
\end{tabular}
\end{table}

\begin{table}[htbp]
\centering
\setlength{\tabcolsep}{4pt}
\caption{Study A, damage and sample efficiency. The uniform-random floor is
$0.26$. Steps to $0.80$ is the median over seeds, with seeds reaching the
threshold in parentheses.}
\small
\begin{tabular}{lccccc}
\toprule
Corruption & Native & Zero-shot & Jumpstart & Steps to 0.80 & Speed-up \\
\midrule
\texttt{large\_multigoal} & 0.229 & 0.274 & $+0.059$ & 84,000 (27/30) & $1.80\times$ \\
\texttt{action\_gain\_jitter} & 0.818 & 0.821 & $+0.641$ & 0 (30/30) & at step 0 \\
\texttt{negate\_both} (ref.) & 0.825 & 0.000 & $-0.217$ & 151,200 (27/30) & $1.00\times$ \\
\midrule
GFlowNet scratch & --- & --- & --- & 151,200 (28/30) & $1.00\times$ \\
\bottomrule
\end{tabular}
\end{table}

\paragraph{Findings.}
\begin{itemize}
  \item \textbf{Success rate never improves, in any condition.} All three differences are negative ($-0.007$, $-0.021$, $-0.025$) and none is significant. Pretraining does not raise the converged success rate.
  \item \textbf{Route quality does improve, consistently.} Steps-to-goal falls in all three conditions, significantly so for \texttt{large\_multigoal} ($-11.5$, $g=-0.92$, $p_{\mathrm{Holm}}=0.003$), and path efficiency rises in all three. The PPO initialisation transfers a directness the GFlowNet does not otherwise acquire within budget.
  \item \textbf{Fine-tuning can destroy a good policy.} Under \texttt{action\_gain\_jitter} the pretrained PPO policy enters adaptation at $0.858$ success but ends at $0.704$, \emph{below} the $0.726$ scratch control. The GFlowNet objective samples trajectories in proportion to reward rather than maximising it, so continued training pulls an already near-optimal policy back toward its sampling distribution. A jumpstart of $+0.641$ is therefore entirely surrendered.
  \item \textbf{The ceiling is set by the fine-tuning algorithm, not the initialisation.} Every arm converges near $0.70$--$0.73$, irrespective of where it started.
\end{itemize}

\clearpage
\section{Study B: GFlowNet to PPO}

The same design with the algorithms exchanged: a GFlowNet pretrains on the
corrupted environment and PPO is fine-tuned on the original, against a PPO
from-scratch control.

\begin{table}[htbp]
\centering
\setlength{\tabcolsep}{4pt}
\caption{Study B primary contrast: fine-tuned $-$ scratch on the original
environment. Both arms are PPO. Conventions as in Study A;
${}^{*}p<0.05$, ${}^{***}p<0.001$, Holm-corrected.}
\small
\begin{tabular}{lcccc}
\toprule
Corruption & Success rate & Mean return & Steps to goal & Path efficiency \\
\midrule
\texttt{large\_multigoal} & $-0.002$ & $-7.3$ & $+6.3^{*}$ & $-0.020^{*}$ \\
 & \footnotesize$(-0.04)$ & \footnotesize$(-0.53)$ & \footnotesize$(0.80)$ & \footnotesize$(-0.72)$ \\[3pt]
\texttt{action\_gain\_jitter} & $+0.011$ & $+8.2^{*}$ & $-2.6$ & $-0.002$ \\
 & \footnotesize$(0.35)$ & \footnotesize$(0.71)$ & \footnotesize$(-0.46)$ & \footnotesize$(-0.08)$ \\[3pt]
\texttt{negate\_both} (ref.) & $-0.051^{***}$ & $-40.1^{***}$ & $+3.8$ & $-0.076^{***}$ \\
 & \footnotesize$(-1.00)$ & \footnotesize$(-1.61)$ & \footnotesize$(0.45)$ & \footnotesize$(-1.87)$ \\
\bottomrule
\end{tabular}
\end{table}

\begin{table}[htbp]
\centering
\setlength{\tabcolsep}{4pt}
\caption{Study B, final performance by arm, mean $\pm$ SD over $n=30$ seeds.
Conventions as in Study A. The PPO scratch arm is shared across both
corruptions.}
\small
\begin{tabular}{llcccc}
\toprule
Corruption & Arm & Success & Return & Steps & Path eff. \\
\midrule
\multicolumn{6}{l}{\texttt{large\_multigoal}} \\
 & Pretrain (native)    & $0.178 \pm 0.080$ & $8.6 \pm 6.1$ & $154.5 \pm 37.7$ & $0.788 \pm 0.085$ \\
 & Pretrain (zero-shot) & $0.311 \pm 0.142$ & $16.9 \pm 6.9$ & $84.8 \pm 21.4$ & $0.666 \pm 0.081$ \\
 & Fine-tuned           & $0.825 \pm 0.046$ & $178.4 \pm 13.4$ & $74.9 \pm 9.0$ & $0.879 \pm 0.030$ \\
\midrule
\multicolumn{6}{l}{\texttt{action\_gain\_jitter}} \\
 & Pretrain (native)    & $0.708 \pm 0.182$ & $48.5 \pm 20.6$ & $83.9 \pm 15.5$ & $0.673 \pm 0.049$ \\
 & Pretrain (zero-shot) & $0.703 \pm 0.183$ & $48.6 \pm 20.7$ & $83.5 \pm 15.2$ & $0.675 \pm 0.048$ \\
 & Fine-tuned           & $0.838 \pm 0.035$ & $193.9 \pm 8.7$ & $66.0 \pm 5.0$ & $0.897 \pm 0.030$ \\
\midrule
\multicolumn{2}{l}{PPO scratch (control)} & $0.827 \pm 0.029$ & $185.7 \pm 13.7$ & $68.6 \pm 6.2$ & $0.899 \pm 0.024$ \\
\bottomrule
\end{tabular}
\end{table}

\begin{table}[htbp]
\centering
\setlength{\tabcolsep}{4pt}
\caption{Study B, damage and sample efficiency. Conventions as in Study A.}
\small
\begin{tabular}{lccccc}
\toprule
Corruption & Native & Zero-shot & Jumpstart & Steps to 0.80 & Speed-up \\
\midrule
\texttt{large\_multigoal} & 0.178 & 0.311 & $+0.042$ & 150,000 (29/30) & $1.25\times$ \\
\texttt{action\_gain\_jitter} & 0.708 & 0.703 & $+0.467$ & 30,000 (30/30) & $6.25\times$ \\
\texttt{negate\_both} (ref.) & 0.666 & 0.017 & $-0.220$ & 232,500 (22/30) & $0.81\times$ \\
\midrule
PPO scratch & --- & --- & --- & 187,500 (30/30) & $1.00\times$ \\
\bottomrule
\end{tabular}
\end{table}

\paragraph{Findings.}
\begin{itemize}
  \item \textbf{Corruption severity spans the full range.} \texttt{negate\_both} leaves a policy that is actively wrong ($0.017$, below the random floor); \texttt{large\_multigoal} mediocre but not wrong ($0.311$); \texttt{action\_gain\_jitter} undamaged ($0.703$ vs.\ $0.708$ native, $p=0.92$).
  \item \textbf{No corruption improved final performance.} Success rate differences are $-0.002$ ($p_{\mathrm{Holm}}=0.868$) and $+0.011$ ($p_{\mathrm{Holm}}=0.354$); the reference condition is significantly worse. As in Study A, pretraining is neutral to harmful at convergence, never beneficial.
  \item \textbf{The benefit is sample efficiency only.} \texttt{action\_gain\_jitter} gives a $+0.467$ jumpstart and reaches $0.80$ success $6.25\times$ sooner, then converges to a statistical tie.
  \item \textbf{Structural corruption leaves a residual cost.} \texttt{large\_multigoal} matches on success but is $6.3$ steps slower ($g=0.80$) and $0.020$ less path-efficient ($g=-0.72$), both significant after correction --- the opposite sign to Study A, where the PPO initialisation \emph{shortened} routes.
  \item \textbf{Initialisation washes out.} Across-seed SD falls from $\pm0.182$ (pretrained) to $\pm0.035$ (fine-tuned), against $\pm0.029$ for the control: PPO converges to a common ceiling regardless of where it starts.
\end{itemize}

\clearpage
\section{Cross-Study Comparison}

\begin{table}[htbp]
\centering
\setlength{\tabcolsep}{4pt}
\caption{Both directions, \texttt{action\_gain\_jitter}. The entry point is the
pretrained policy's zero-shot score; the exit point is the fine-tuned arm at
300,000 steps.}
\small
\begin{tabular}{lccc}
\toprule
 & Entry (zero-shot) & Exit (fine-tuned) & Scratch control \\
\midrule
Study A (PPO $\rightarrow$ GFN) & 0.821 & $0.704$ & $0.726$ \\
Study B (GFN $\rightarrow$ PPO) & 0.703 & $0.838$ & $0.827$ \\
\bottomrule
\end{tabular}
\end{table}

\begin{itemize}
  \item \textbf{Neither direction raises the converged success rate.} Across six corruption--direction pairs, not one shows a significant gain on the primary metric. The consistent result is that cross-algorithm pretraining does not improve final task performance in this setting.
  \item \textbf{The two directions fail differently.} Study A converts its head start into route quality --- shorter, more direct paths --- while success rate is flat or slightly negative. Study B converts its head start into sample efficiency ($6.25\times$) and then ties.
  \item \textbf{Direction determines the ceiling.} GFlowNet fine-tuning converges near $0.72$ and PPO fine-tuning near $0.83$, regardless of initialisation quality. In Study A under \texttt{action\_gain\_jitter} this is strong enough to reverse the transfer: a policy entering at $0.821$ exits at $0.704$.
\end{itemize}

\section{Limitations}

\begin{itemize}
  \item \texttt{large\_multigoal} is not difficulty-preserving --- maze, size, start and goal distributions change together --- so its zero-shot number conflates a wrong policy with an unseen task shape, and its native score is not comparable across environments. The primary contrasts are unaffected: both arms run the same algorithm on the same environment under the same budget.
  \item Mean return is not independent. Episodes are not terminated on success and accrue reward while at the goal, so return $\approx$ success rate $\times$ remaining time ($0.838\times(300-66)\approx196$ vs.\ measured $193.9$). It aggregates success and latency rather than measuring a distinct property, which is why it can be the only significant metric while success rate is flat.
  \item At $n=30$, power is $\approx0.93$ for $|g|=1.1$ but $0.3$--$0.5$ for $|g|\approx0.5$. Non-significant entries should be read as unresolved, not as established nulls; effect sizes are therefore reported unconditionally.
  \item Single environment family, one action interface, one 300,000-step budget. Results do not extend to continuous-action settings, dense rewards, or longer budgets without re-measurement.
\end{itemize}
